% =====================================================================
\chapter{Soluciones}
\label{ap:soluciones}
% =====================================================================

Los resultados finales se dan con tres cifras significativas y los intermedios con cuatro, para que el redondeo no se acumule. Las remisiones «M16, §4.3» indican el módulo y la sección de sus apuntes. Al final de cada solución figuran, entre corchetes, los criterios de corrección.

% ---------------------------------------------------------------------
\section{Parte A: clave de respuestas}

\begin{table}[h]
\centering\footnotesize
\caption{Clave de la Parte~A.}\label{tab:clave}
\begin{tabularx}{\textwidth}{@{}c c >{\raggedright\arraybackslash}X l@{}}
\toprule
N.º & Resp. & Justificación & Repasar \\
\midrule
1 & (c) & $6\cdot 3\cdot 10^{9}\cdot 6\cdot 10^{10}=1{,}08\cdot 10^{21}$ FLOP. & M17, §1.2 \\
2 & (a) & $20\cdot 8\cdot 10^{9}=1{,}6\cdot 10^{11}$ tokens. & M16, §4.5 \\
3 & (d) & Chinchilla escala $N$ y $D$ casi por igual ($a\approx 0{,}5$); Kaplan daba $a=0{,}73$. & M16, §5.2 \\
4 & (b) & ZeRO-3: $16\Psi/n_{\mathrm{d}}=16\cdot 1{,}3\cdot 10^{10}/16=1{,}3\cdot 10^{10}$ bytes. & M17, §5.6 \\
5 & (c) & $\Icrit=989\cdot 10^{12}/(3{,}35\cdot 10^{12})=295{,}2$; 156 es la A100. & M17, §2.2 \\
6 & (a) & $\Bmem/(sN)=3{,}35\cdot 10^{12}/(6\cdot 10^{10})=55{,}8$ tokens/s; los FLOP no intervienen. & M19, §2.3 \\
7 & (d) & $2\cdot 40\cdot 8\cdot 128\cdot 2=163\,840$ bytes $=160$~KiB. & M19, §2.2 \\
8 & (b) & $(p_{\mathrm{e}}-1)/(m+p_{\mathrm{e}}-1)=7/31=0{,}2258$. & M17, §5.5 \\
9 & (a) & $0{,}25\cdot 3{,}14\cdot 10^{11}=7{,}85\cdot 10^{10}$. & M18, §4.6 \\
10 & (c) & Cada token sigue usando 2 expertos (mismos FLOP), pero hay que guardar los 64. & M18, §4.2 \\
11 & (d) & Las cabezas de consulta comparten $K$ y $V$: la caché se divide por $h/\nkv$. & M18, §2.3 \\
12 & (b) & $24\cdot 2048=49\,152$ (el valor exacto es $24\cdot 2047+1=49\,129$). & M18, §3.2 \\
13 & (a) & Con $P_i=1/4$: $4\cdot\frac14\sum_i f_i=1$, aunque $f$ esté desequilibrado. & M18, §4.3 \\
14 & (c) & Resultado exacto con $O(n^{2}d)$ FLOP y $O(n)$ memoria adicional: no escribe $S$ en HBM. & M17, §3.4 \\
15 & (d) & Procesa todas las posiciones a la vez: productos matriz por matriz, de intensidad alta. & M19, §2.1 \\
16 & (b) & $\min(p,q)+\max(0,p-q)=p$: el muestreo por rechazo modificado es exacto. & M19, §3.2 \\
17 & (a) & $b+b_{\Delta}/g=4+16/64=4{,}25$; $4{,}125$ corresponde a $g=128$. & M19, §4.3 \\
18 & (c) & $0{,}95^{20}=0{,}3585$. & M16, §6.2 \\
19 & (d) & Asigna bloques bajo demanda y elimina la fragmentación; caben lotes mayores. & M19, §6.3 \\
20 & (b) & $\MFU=6N\tau/(n_{\mathrm{d}}\Ppico)$, con $\tau$ los tokens/s (definición de PaLM). & M17, §2.5 \\
\bottomrule
\end{tabularx}
\end{table}

Las cifras y afirmaciones sobre trabajos concretos proceden de sus fuentes públicas: Kaplan \emph{et al.}~\cite{kaplan} y Hoffmann \emph{et al.}~\cite{hoffmann} (cuestiones 2 y 3), ZeRO~\cite{zero} (4), FlashAttention~\cite{flash} (14), el anuncio de Grok-1~\cite{grok1} (9), Leviathan \emph{et al.}~\cite{leviathan} (16), vLLM~\cite{vllm} (19), PaLM~\cite{palm} (20) y Schaeffer \emph{et al.}~\cite{schaeffer} (18). Los apuntes de los módulos~\cite{apuntes} desarrollan cada resultado.

% ---------------------------------------------------------------------
\section{Parte B: soluciones de los problemas}

\begin{solucion}{eje:chinchilla}
\datos $L(N,D)=E+AN^{-\alpha}+BD^{-\beta}$ con $E=1{,}69$, $A=406{,}4$, $B=410{,}7$, $\alpha=0{,}34$, $\beta=0{,}28$; restricción $C=6ND$ con $C$ fijo.
\planteamiento se despeja $D=C/(6N)$, se sustituye en $L$ y se anula la derivada respecto a $N$ (es el teorema de la frontera eficiente, M16, §4.3). Después se evalúa la ley en los repartos pedidos.
\pasos

\noindent\textbf{(a)}
\begin{enumerate}[label=\arabic*.]
  \item Con $c=C/6$, $D=c/N$ y $B/D^{\beta}=Bc^{-\beta}N^{\beta}$, luego $L(N)=E+AN^{-\alpha}+Bc^{-\beta}N^{\beta}$.
  \item $\dfrac{dL}{dN}=-\alpha AN^{-\alpha-1}+\beta Bc^{-\beta}N^{\beta-1}$. Se iguala a~0 y se multiplica por $N>0$: $\alpha AN^{-\alpha}=\beta Bc^{-\beta}N^{\beta}$.
  \item Como $c^{-\beta}N^{\beta}=(c/N)^{-\beta}=D^{-\beta}$, queda $\alpha A/N^{\alpha}=\beta B/D^{\beta}$.
  \item Interpretación: $-N\,\partial L/\partial N=\alpha A/N^{\alpha}$ y $-D\,\partial L/\partial D=\beta B/D^{\beta}$; en el óptimo, aumentar un 1\,\% los parámetros reduce la pérdida lo mismo que aumentar un 1\,\% los datos.
\end{enumerate}

\noindent\textbf{(b)}
\begin{enumerate}[label=\arabic*.]
  \item De la igualdad del paso~2: $N^{\alpha+\beta}=\dfrac{\alpha A}{\beta B}\,c^{\beta}$, luego $N=\Bigl(\dfrac{\alpha A}{\beta B}\Bigr)^{1/(\alpha+\beta)}c^{\beta/(\alpha+\beta)}=G\,c^{a}$.
  \item $\alpha A=0{,}34\cdot 406{,}4=138{,}2$; $\beta B=0{,}28\cdot 410{,}7=115{,}0$; cociente $1{,}2016$.
  \item $a=0{,}28/0{,}62=0{,}4516$. $\ln 1{,}2016=0{,}1836$; $0{,}1836/0{,}62=0{,}2962$; $G=e^{0{,}2962}=1{,}345$.
\end{enumerate}

\noindent\textbf{(c)}
\begin{enumerate}[label=\arabic*.]
  \item $C=10^{24}$: $c=1{,}667\cdot 10^{23}$, $\log_{10}c=23{,}22$; $a\log_{10}c=10{,}49$; $c^{a}=3{,}071\cdot 10^{10}$.
  \item $\Nopt=1{,}345\cdot 3{,}071\cdot 10^{10}=4{,}130\cdot 10^{10}$; $\Dopt=c/\Nopt=4{,}036\cdot 10^{12}$; $\Dopt/\Nopt=97{,}7$.
  \item $C=10^{23}$: $\Nopt$ se divide por $10^{a}=2{,}829$: $\Nopt=1{,}460\cdot 10^{10}$ y $\Dopt=1{,}667\cdot 10^{22}/1{,}460\cdot 10^{10}=1{,}142\cdot 10^{12}$.
  \item Pérdida a $10^{24}$: $N^{0{,}34}=e^{0{,}34\cdot 24{,}44}=4068$, $A/N^{\alpha}=0{,}09990$; $D^{0{,}28}=e^{0{,}28\cdot 29{,}03}=3386$, $B/D^{\beta}=0{,}1213$; $L=1{,}69+0{,}09990+0{,}1213=1{,}911$.
\end{enumerate}

\noindent\textbf{(d)}
\begin{enumerate}[label=\arabic*.]
  \item $C=6N\cdot 20N=120N^{2}$: $N=\sqrt{10^{24}/120}=9{,}129\cdot 10^{10}$; $D=20N=1{,}826\cdot 10^{12}$.
  \item $A/N^{\alpha}=406{,}4/5328=0{,}07628$; $B/D^{\beta}=410{,}7/2712=0{,}1515$; $L=1{,}918$.
  \item Diferencia: $1{,}9177-1{,}9112=0{,}0066$.
\end{enumerate}

\noindent\textbf{(e)}
\begin{enumerate}[label=\arabic*.]
  \item Con $D=2\cdot 10^{12}$: $N=10^{24}/(6\cdot 2\cdot 10^{12})=8{,}333\cdot 10^{10}$.
  \item $A/N^{\alpha}=406{,}4/5165=0{,}07868$; $B/D^{\beta}=410{,}7/2782=0{,}1477$; $L=1{,}916$, es decir, $0{,}0051$ más que el óptimo.
\end{enumerate}
\resultado (a) $\alpha A/N^{\alpha}=\beta B/D^{\beta}$. (b) $a=0{,}452$, $G=1{,}34$. (c) A $10^{23}$~FLOP, $\Nopt\approx 1{,}46\cdot 10^{10}$ y $\Dopt\approx 1{,}14\cdot 10^{12}$; a $10^{24}$~FLOP, $\Nopt\approx 4{,}13\cdot 10^{10}$, $\Dopt\approx 4{,}04\cdot 10^{12}$, $\Dopt/\Nopt\approx 97{,}7$ y $L\approx 1{,}91$. (d) Regla de 20: $N\approx 9{,}13\cdot 10^{10}$, $D\approx 1{,}83\cdot 10^{12}$ y $L\approx 1{,}92$; el óptimo gana por solo $0{,}0066$: el mínimo es muy plano (M16, §4.6). (e) $N\approx 8{,}33\cdot 10^{10}$ y $L\approx 1{,}92$, $0{,}0051$ peor que el óptimo.
\comprobacion en (c), $0{,}34\cdot 0{,}09990=0{,}03397$ y $0{,}28\cdot 0{,}1213=0{,}03396$: se cumple la igualdad de~(a); además, $6\cdot 4{,}130\cdot 10^{10}\cdot 4{,}036\cdot 10^{12}=1{,}000\cdot 10^{24}$. El término de datos es $0{,}1213/0{,}09990=1{,}214=\alpha/\beta$ veces el de parámetros, como dice el corolario de equilibrio (M16, §4.3). En~(e), el punto queda entre el óptimo y la regla de 20, y su pérdida también: $1{,}911<1{,}916<1{,}918$. La función $\phi(u)=Ae^{-\alpha u}+Bc^{-\beta}e^{\beta u}$, con $u=\ln N$, es estrictamente convexa, así que el punto crítico es el mínimo.

[Criterios: (a) 1~p por sustituir la restricción, 1~p por la derivada y 1~p por la interpretación; se admite la resolución con multiplicadores de Lagrange. (b) 2~p. (c) 1~p por cada presupuesto. (d) 2~p, o 1~p si solo se calculan las cifras. (e) 1~p.]

Repasar: M16, §4.2, §4.3, §4.5 y §4.6.
\end{solucion}

% ---------------------------------------------------------------------
\begin{solucion}{eje:fermi}
\datos $N=7\cdot 10^{10}$; $D=1{,}5\cdot 10^{13}$ tokens; 8192 H100 a 989~TFLOP/s; $\MFU=0{,}4$; 700~W; $\PUE=1{,}2$; 2{,}5~\$/GPU-hora. GPT-2 XL: $V=50\,257$, $L=48$, $d=1600$, $\dff=4288$; 495~TFLOP/s en TF32 con $\MFU=0{,}5$.
\planteamiento $C\approx 6ND$ (M17, §1.2); duración $t=C/(n_{\mathrm{d}}\Ppico\MFU)$ (M17, §6.1); energía $\approx$ GPU-horas $\cdot P_{\mathrm{GPU}}\cdot\PUE$ (M17, §6.2); coste de inferencia $\approx 2N$ FLOP por token (M16, §5.3).
\pasos

\noindent\textbf{(a)}
\begin{enumerate}[label=\arabic*.]
  \item $C=6\cdot 7\cdot 10^{10}\cdot 1{,}5\cdot 10^{13}=6{,}3\cdot 10^{24}$~FLOP.
\end{enumerate}

\noindent\textbf{(b)}
\begin{enumerate}[label=\arabic*.]
  \item Rendimiento útil por GPU: $0{,}4\cdot 989\cdot 10^{12}=3{,}956\cdot 10^{14}$~FLOP/s.
  \item GPU-segundos: $6{,}3\cdot 10^{24}/3{,}956\cdot 10^{14}=1{,}593\cdot 10^{10}$~s; GPU-horas: $1{,}593\cdot 10^{10}/3600=4{,}424\cdot 10^{6}$.
  \item Horas de reloj: $4{,}424\cdot 10^{6}/8192=540{,}0$~h, es decir, $22{,}50$ días.
  \item Energía: $4{,}424\cdot 10^{6}~\text{h}\cdot 0{,}7~\text{kW}\cdot 1{,}2=3{,}716\cdot 10^{6}$~kWh $=3716$~MWh.
  \item Coste: $4{,}424\cdot 10^{6}\cdot 2{,}5=1{,}106\cdot 10^{7}$~\$.
\end{enumerate}

\noindent\textbf{(c)}
\begin{enumerate}[label=\arabic*.]
  \item \emph{Embeddings} no compartidos: $2Vd=2\cdot 50\,257\cdot 1600=160\,822\,400$.
  \item Por capa: atención $4d^{2}=10\,240\,000$; MLP SwiGLU $3d\dff=3\cdot 1600\cdot 4288=20\,582\,400$; dos RMSNorm, $2d=3200$. Suma: $30\,825\,600$.
  \item Total: $48\cdot 30\,825\,600+160\,822\,400+1600=1\,640\,452\,800\approx 1{,}640\cdot 10^{9}$ (el último sumando es la RMSNorm final). En fp32: $4\cdot 1{,}640\cdot 10^{9}=6{,}562\cdot 10^{9}$ bytes.
  \item Tokens: $4\cdot 10^{5}\cdot 1024\cdot 1024=4{,}194\cdot 10^{11}$. Cómputo: $6\cdot 1{,}640\cdot 10^{9}\cdot 4{,}194\cdot 10^{11}=4{,}128\cdot 10^{21}$~FLOP.
  \item $t=4{,}128\cdot 10^{21}/(0{,}5\cdot 495\cdot 10^{12})=1{,}668\cdot 10^{7}$~s $=4633$~h $=193{,}1$ días.
\end{enumerate}

\noindent\textbf{(d)}
\begin{enumerate}[label=\arabic*.]
  \item GPU-segundos declarados: $6{,}4\cdot 10^{6}\cdot 3600=2{,}304\cdot 10^{10}$~s; FLOP de pico disponibles: $2{,}304\cdot 10^{10}\cdot 989\cdot 10^{12}=2{,}279\cdot 10^{25}$.
  \item $\MFU_{\mathrm{ef}}=6{,}3\cdot 10^{24}/2{,}279\cdot 10^{25}=0{,}2765$.
\end{enumerate}

\noindent\textbf{(e)}
\begin{enumerate}[label=\arabic*.]
  \item $D/N=1{,}5\cdot 10^{13}/7\cdot 10^{10}=214{,}3$ tokens por parámetro, $10{,}71$ veces la regla.
  \item Regla de 20 con el mismo $C$: $N=\sqrt{6{,}3\cdot 10^{24}/120}=2{,}291\cdot 10^{11}$ y $D=20N=4{,}583\cdot 10^{12}$.
  \item Servir $10^{13}$ tokens: $2\cdot 7\cdot 10^{10}\cdot 10^{13}=1{,}4\cdot 10^{24}$~FLOP con el modelo de 70\,000 millones, frente a $2\cdot 2{,}291\cdot 10^{11}\cdot 10^{13}=4{,}583\cdot 10^{24}$ con el de la regla.
\end{enumerate}
\resultado (a) $C\approx 6{,}30\cdot 10^{24}$~FLOP. (b) $4{,}42\cdot 10^{6}$ GPU-horas, $22{,}5$ días, $\approx 3720$~MWh y $\approx 1{,}11\cdot 10^{7}$~\$. (c) $N\approx 1{,}64\cdot 10^{9}$ parámetros, $6{,}56$~GB en fp32 y $\approx 4630$~h, unos 193 días en una sola GPU. (d) $\MFU_{\mathrm{ef}}\approx 27{,}6\,\%$; razones: la ficha cuenta todo el tiempo de GPU (reinicios tras fallos, puntos de control, evaluaciones y experimentos), la comunicación no solapada, la burbuja del pipeline y los FLOP de la atención y de la fase de contexto largo, que $6ND$ no cuenta. (e) $\approx 214$ tokens por parámetro, unas 11 veces la regla; la regla daría $N\approx 2{,}29\cdot 10^{11}$ y $D\approx 4{,}58\cdot 10^{12}$. Servir $10^{13}$ tokens cuesta $1{,}40\cdot 10^{24}$~FLOP con el modelo pequeño y $4{,}58\cdot 10^{24}$ con el grande: con un uso intenso, el coste de inferencia es del orden del de entrenamiento, y un modelo más pequeño y más entrenado resulta más barato en su ciclo de vida.
\comprobacion (b) $8192\cdot 22{,}5\cdot 24=4{,}424\cdot 10^{6}$ GPU-horas. (c) Con la atención, cada token añade unos $4Ld\,T$ FLOP hacia delante con $T=1024$; frente a $2N$, son $4\cdot 48\cdot 1600\cdot 1024/(2\cdot 1{,}640\cdot 10^{9})=9{,}6\,\%$ más, de modo que la cifra de 193 días es una cota inferior (ampliación: la tarea de CS336 pide contar también esos productos). (d) $4{,}424\cdot 10^{6}/6{,}4\cdot 10^{6}=0{,}6913$ y $0{,}6913\cdot 40\,\%=27{,}65\,\%$. (e) $4{,}583/1{,}4=3{,}27=2{,}291\cdot 10^{11}/7\cdot 10^{10}$: la inferencia escala con $N$.

[Criterios: (a) 1~p. (b) 0{,}75~p por cada cifra. (c) 1~p por $N$ y la memoria, 1~p por el tiempo. (d) 1~p por la MFU y 1~p por las razones. (e) 1~p por las cifras y 1~p por el argumento.]

Repasar: M17, §1.2, §2.5, §6.1 y §6.2; M16, §5.3.
\end{solucion}

% ---------------------------------------------------------------------
\begin{solucion}{eje:paralelismo}
\datos $\Psi=7\cdot 10^{10}$; $\mathrm{TP}=8$, $\mathrm{PP}=p_{\mathrm{e}}=4$, $\mathrm{DP}=64$; 2048 H100; lote de 1024 secuencias de 4096 tokens; microlote de una secuencia; NVLink $\approx 900$~GB/s e InfiniBand $\approx 50$~GB/s por GPU.
\planteamiento $n_{\mathrm{d}}=\mathrm{TP}\cdot\mathrm{PP}\cdot\mathrm{DP}$. Memoria de estado con la tabla de ZeRO (M17, §5.6). Burbuja $(p_{\mathrm{e}}-1)/(m+p_{\mathrm{e}}-1)$ (M17, §5.5). All-reduce en anillo: $t_{\mathrm{AR}}=2(n-1)S/(nW)$ (M17, §5.1).
\pasos

\noindent\textbf{(a)}
\begin{enumerate}[label=\arabic*.]
  \item $8\cdot 4\cdot 64=2048=256\cdot 8$: cuadra.
  \item TP comunica activaciones varias veces por capa y necesita el enlace más rápido. Con $\mathrm{TP}=16$, cada grupo ocuparía dos nodos y sus all-reduce cruzarían InfiniBand, $900/50=18$ veces más lento que NVLink.
\end{enumerate}

\noindent\textbf{(b)}
\begin{enumerate}[label=\arabic*.]
  \item $16\Psi=1{,}12\cdot 10^{12}$~bytes; entre $\mathrm{TP}\cdot\mathrm{PP}=32$ GPU: $3{,}5\cdot 10^{10}$~bytes $=35$~GB por GPU. Quedan $80-35=45$~GB para activaciones.
  \item Paralelismo de datos puro: cada GPU guardaría los $1{,}12\cdot 10^{12}$~bytes $=1120$~GB, 14 veces su memoria. Ni siquiera los pesos en bf16 ($140$~GB) caben. No es viable.
\end{enumerate}

\noindent\textbf{(c)}
\begin{enumerate}[label=\arabic*.]
  \item Pesos y gradientes: $4\Psi/32=2{,}8\cdot 10^{11}/32=8{,}75\cdot 10^{9}$~bytes.
  \item Optimizador: $12\Psi/(32\cdot 64)=8{,}4\cdot 10^{11}/2048=4{,}102\cdot 10^{8}$~bytes.
  \item Total: $8{,}75+0{,}4102=9{,}160$~GB por GPU; se recuperan $35-9{,}160=25{,}84$~GB.
\end{enumerate}

\noindent\textbf{(d)}
\begin{enumerate}[label=\arabic*.]
  \item Tokens por paso: $1024\cdot 4096=4\,194\,304$.
  \item Secuencias por réplica: $1024/64=16$; con microlotes de una secuencia, $m=16$.
  \item Burbuja: $(4-1)/(16+4-1)=3/19=0{,}1579$.
\end{enumerate}

\noindent\textbf{(e)}
\begin{enumerate}[label=\arabic*.]
  \item FLOP por paso: $6\cdot 7\cdot 10^{10}\cdot 4{,}194\cdot 10^{6}=1{,}762\cdot 10^{18}$. Rendimiento útil del clúster: $2048\cdot 989\cdot 10^{12}\cdot 0{,}4=8{,}102\cdot 10^{17}$~FLOP/s. $t_{\mathrm{c}}=2{,}174$~s.
  \item Gradientes bf16 por GPU: $S=2\Psi/32=4{,}375\cdot 10^{9}$~bytes. Volumen enviado: $2\cdot(63/64)\cdot 4{,}375\cdot 10^{9}=8{,}613\cdot 10^{9}$~bytes. $t_{\mathrm{AR}}=8{,}613\cdot 10^{9}/(50\cdot 10^{9})=0{,}1723$~s, el $7{,}9\,\%$ de $t_{\mathrm{c}}$.
  \item El cómputo por secuencia y réplica es $t_{\mathrm{c}}/16=0{,}1359$~s, mientras que $t_{\mathrm{AR}}$ no depende del lote. Con $s$ secuencias por réplica, la comunicación queda por debajo del cómputo si $0{,}1359\,s\ge 0{,}1723$, es decir, $s\ge 1{,}268$: hacen falta al menos 2 secuencias por réplica.
\end{enumerate}

\noindent\textbf{(f)}
\begin{enumerate}[label=\arabic*.]
  \item Pasos: $1{,}5\cdot 10^{13}/4{,}194\cdot 10^{6}=3{,}577\cdot 10^{6}$; tiempo: $3{,}577\cdot 10^{6}\cdot 2{,}174=7{,}776\cdot 10^{6}$~s $=90{,}00$ días.
\end{enumerate}
\resultado (a) 2048 GPU; $\mathrm{TP}=8$ cabe en un nodo y usa NVLink. (b) $35{,}0$~GB por GPU: cabe; el paralelismo de datos puro exigiría 1120~GB por GPU y no es viable. (c) $\approx 9{,}16$~GB por GPU. (d) $4\,194\,304$ tokens por paso, $m=16$ y burbuja $\approx 15{,}8\,\%$. (e) $t_{\mathrm{c}}\approx 2{,}17$~s y $t_{\mathrm{AR}}\approx 0{,}172$~s ($7{,}9\,\%$); como además se solapa con la pasada hacia atrás, no es el cuello de botella: pesa más la burbuja. Con una sola secuencia por réplica sí lo sería; bastan 2. (f) $\approx 3{,}58\cdot 10^{6}$ pasos y $\approx 90{,}0$ días.
\comprobacion (c) El optimizador era el $12/16=75\,\%$ de los 35~GB, $26{,}25$~GB; ahora ocupa $0{,}41$~GB, y $26{,}25-0{,}41=25{,}84$~GB. (d) Con $m\to\infty$ la burbuja tiende a~0 y con $m=1$ vale $3/4$. (e) Para que la comunicación igualase al cómputo con el lote completo, el enlace debería bajar a $8{,}613/2{,}174=3{,}96$~GB/s. (f) Son 4 veces los 22{,}5 días del \cref{eje:fermi}, porque hay $8192/2048=4$ veces menos GPU con la misma MFU.

[Criterios: (a) 1~p. (b) 1~p por la memoria y 1~p por el paralelismo de datos puro. (c) 2~p. (d) 2~p. (e) 1~p por los dos tiempos y 1~p por el lote mínimo. (f) 1~p.]

Repasar: M17, §5.1, §5.3, §5.4, §5.5 y §5.6.
\end{solucion}

% ---------------------------------------------------------------------
\begin{solucion}{eje:moe}
\datos $d=5120$; $L=40$; 16 expertos top-2; $\dff=8192$, SwiGLU con 3 matrices; 40 cabezas de consulta, $\nkv=8$, $\dhd=128$; $V=64\,000$ sin compartir.
\planteamiento se suman expertos, atención y \emph{embeddings}; los activos cuentan 2 expertos por capa (M18, §4.2). FLOP hacia delante $\approx 2\Nact$ por token; memoria $=2\Ntot$ bytes en bf16. Pesos del enrutador como softmax entre los elegidos (M18, §4.1); $\Lbal$ (M18, §4.3); capacidad $C_{\mathrm{exp}}=c\,kB_{\mathrm{t}}/N$ (M18, §4.4).
\pasos

\noindent\textbf{(a)}
\begin{enumerate}[label=\arabic*.]
  \item Experto: $3\cdot 5120\cdot 8192=125\,829\,120$.
  \item Atención: $K$ y $V$ tienen dimensión $8\cdot 128=1024$. $W_Q$ y $W_O$: $2\cdot 5120^{2}=52\,428\,800$; $W_K$ y $W_V$: $2\cdot 5120\cdot 1024=10\,485\,760$. Suma: $62\,914\,560$.
  \item \emph{Embeddings}: $2\cdot 64\,000\cdot 5120=655\,360\,000$.
  \item Total: $40\cdot(16\cdot 125\,829\,120+62\,914\,560)+655\,360\,000=40\cdot 2\,076\,180\,480+655\,360\,000=83\,702\,579\,200$.
  \item Activos: $40\cdot(2\cdot 125\,829\,120+62\,914\,560)+655\,360\,000=40\cdot 314\,572\,800+655\,360\,000=13\,238\,272\,000$.
  \item Fracción activa: $13\,238\,272\,000/83\,702\,579\,200=0{,}1582$.
\end{enumerate}

\noindent\textbf{(b)}
\begin{enumerate}[label=\arabic*.]
  \item FLOP por token: $2\cdot 1{,}324\cdot 10^{10}=2{,}648\cdot 10^{10}$.
  \item Memoria: $2\cdot 8{,}370\cdot 10^{10}=1{,}674\cdot 10^{11}$~bytes $=167{,}4$~GB; $167{,}4/80=2{,}093$, luego hacen falta al menos 3 H100.
\end{enumerate}

\noindent\textbf{(c)}
\begin{enumerate}[label=\arabic*.]
  \item Los dos logits mayores son $2{,}6$ (experto~6) y $1{,}9$ (experto~2); el siguiente, $1{,}7$ (experto~13), queda fuera.
  \item Con dos elegidos, $g_6=1/(1+e^{-\Delta})$ con $\Delta=2{,}6-1{,}9=0{,}7$: $e^{-0{,}7}=0{,}4966$ y $g_6=1/1{,}4966=0{,}6682$; $g_2=1-0{,}6682=0{,}3318$.
\end{enumerate}

\noindent\textbf{(d)}
\begin{enumerate}[label=\arabic*.]
  \item $\sum f_i=4\cdot 0{,}16+12\cdot 0{,}03=0{,}64+0{,}36=1$; $\sum P_i=4\cdot 0{,}12+12\cdot 13/300=0{,}48+0{,}52=1$.
  \item $\sum f_iP_i=4\cdot 0{,}16\cdot 0{,}12+12\cdot 0{,}03\cdot 13/300=0{,}0768+0{,}0156=0{,}0924$.
  \item $\Lbal=16\cdot 0{,}0924=1{,}478$.
\end{enumerate}

\noindent\textbf{(e)}
\begin{enumerate}[label=\arabic*.]
  \item Asignaciones: $2\cdot 8192=16\,384$. Capacidad: $1{,}25\cdot 16\,384/16=1280$ plazas.
  \item Asignaciones a un experto saturado: $0{,}16\cdot 16\,384=2621$; descartes: $2621-1280=1341$.
\end{enumerate}
\resultado (a) Experto, $1{,}26\cdot 10^{8}$; atención por capa, $6{,}29\cdot 10^{7}$; \emph{embeddings}, $6{,}55\cdot 10^{8}$; total $\approx 8{,}37\cdot 10^{10}$; activos $\approx 1{,}32\cdot 10^{10}$; fracción activa $\approx 15{,}8\,\%$. El enrutador añade solo $5120\cdot 16\cdot 40=3\,276\,800$ parámetros. (b) $\approx 2{,}65\cdot 10^{10}$~FLOP por token, como un denso de $1{,}3\cdot 10^{10}$; 167~GB de pesos y al menos 3 H100. Una MoE ahorra cómputo por token, pero no memoria de pesos ni caché KV. (c) Expertos 6 y 2, con $g_6\approx 0{,}668$ y $g_2\approx 0{,}332$. (d) $\Lbal\approx 1{,}48$, frente a~1 con reparto perfecto; el gradiente $\partial\Lbal/\partial P_i=Nf_i$ es mayor en los expertos saturados y reduce su probabilidad. (e) 1280 plazas por experto; cada experto saturado descarta $\approx 1341$ asignaciones, que en esa capa solo pasan por la conexión residual (o por el otro experto elegido).
\comprobacion (a) $\Ntot-\Nact=70\,464\,307\,200=40\cdot 14\cdot 125\,829\,120$: los 14 expertos no usados de cada capa. (b) La caché KV solo depende de $L$, $\nkv$ y $\dhd$, iguales que en un denso con la misma atención. (c) $e^{2{,}6}=13{,}46$ y $e^{1{,}9}=6{,}686$; $13{,}46/(13{,}46+6{,}686)=0{,}6682$. (d) Reparto perfecto: $16\cdot 16\cdot\frac{1}{16}\cdot\frac{1}{16}=1$; el valor $1{,}478$ está en $[1,16]$. (e) Se pierde $1341/2621=51\,\%$ del tráfico de esos expertos: un desequilibrio de $0{,}16$ frente al ideal $1/16=0{,}0625$ ($2{,}56$ veces) no cabe con un factor $1{,}25$.

[Criterios: 2~p por apartado. En (c), 1~p si se aplica la softmax a los 16 logits sin renormalizar.]

Repasar: M18, §4.1, §4.2, §4.3 y §4.4.
\end{solucion}

% ---------------------------------------------------------------------
\begin{solucion}{eje:servicio}
\datos $L=56$; $d=7168$; $\nkv=8$; $\dhd=128$; $N=3{,}5\cdot 10^{10}$; bf16 ($s=2$); 2 H100: 160~GB, $\Bmem=6{,}7$~TB/s y $1{,}978\cdot 10^{15}$~FLOP/s; 10~GB reservados.
\planteamiento caché por token $\ckv=2L\nkv\dhd s$ (M19, §2.2). En la decodificación, $t_{\mathrm{paso}}\approx(\text{bytes de pesos}+\text{bytes de caché})/\Bmem$ (M19, §2.3--§2.4). El \emph{prefill} está limitado por el cómputo (M19, §2.1). Coste de las escalas: $b+b_{\Delta}/g$ bits por peso (M19, §4.3). Decodificación especulativa: $\E[N_{\mathrm{r}}]=(1-\alpha^{\gamma+1})/(1-\alpha)$ tokens por ronda con un coste de $\gamma c+1$ pasadas (M19, §3.4).
\pasos

\noindent\textbf{(a)}
\begin{enumerate}[label=\arabic*.]
  \item $12\cdot 56\cdot 7168^{2}=672\cdot 51\,380\,224=3{,}453\cdot 10^{10}\approx 3{,}5\cdot 10^{10}$.
  \item Pesos: $2\cdot 3{,}5\cdot 10^{10}=7\cdot 10^{10}$~bytes $=70$~GB. Libre: $160-70-10=80$~GB.
  \item $\ckv=2\cdot 56\cdot 8\cdot 128\cdot 2=229\,376$~bytes $=224$~KiB. Por secuencia: $229\,376\cdot 8192=1{,}879\cdot 10^{9}$~bytes $=1{,}75$~GiB.
\end{enumerate}

\noindent\textbf{(b)}
\begin{enumerate}[label=\arabic*.]
  \item GQA: $80/1{,}879=42{,}58$, luego caben 42 secuencias.
  \item MHA: $2\cdot 56\cdot 56\cdot 128\cdot 2=1\,605\,632$~bytes por token (7 veces más); por secuencia, $1{,}315\cdot 10^{10}$~bytes; $80/13{,}15=6{,}08$, luego caben 6.
\end{enumerate}

\noindent\textbf{(c)}
\begin{enumerate}[label=\arabic*.]
  \item (i) Bytes por paso: $70+1{,}879=71{,}88$~GB; $t=71{,}88\cdot 10^{9}/6{,}7\cdot 10^{12}=10{,}73$~ms; $1/t=93{,}21$ tokens/s.
  \item (ii) Caché: $32\cdot 4096\cdot 229\,376=3{,}006\cdot 10^{10}$~bytes; $t=(70+30{,}06)\cdot 10^{9}/6{,}7\cdot 10^{12}=14{,}94$~ms; totales $32/0{,}01494=2142$ tokens/s; por usuario, $66{,}96$ tokens/s.
  \item Cómputo en (ii): $2\cdot 3{,}5\cdot 10^{10}\cdot 32=2{,}24\cdot 10^{12}$~FLOP; $2{,}24\cdot 10^{12}/1{,}978\cdot 10^{15}=1{,}132$~ms $<14{,}94$~ms.
  \item Lote crítico: un paso con lote $B$ hace $2NB$~FLOP y lee $2N$~bytes de pesos, con intensidad $I=B$~FLOP/byte; pasa a estar limitado por el cómputo cuando $B\ge\Icrit=989/3{,}35=295{,}2$ (el cociente es el mismo con una o con dos GPU).
  \item Memoria de ese lote: $295\cdot 4096\cdot 229\,376=2{,}772\cdot 10^{11}$~bytes $=277{,}2$~GB $>80$~GB; con 4096 tokens solo caben $80/0{,}9395=85$ secuencias.
\end{enumerate}

\noindent\textbf{(d)}
\begin{enumerate}[label=\arabic*.]
  \item FLOP: $2\cdot 3{,}5\cdot 10^{10}\cdot 4000=2{,}8\cdot 10^{14}$; rendimiento útil: $0{,}5\cdot 1{,}978\cdot 10^{15}=9{,}89\cdot 10^{14}$~FLOP/s; $\TTFT=0{,}2831$~s.
\end{enumerate}

\noindent\textbf{(e)}
\begin{enumerate}[label=\arabic*.]
  \item Bits por peso: $4+16/64=4{,}25$; pesos: $3{,}5\cdot 10^{10}\cdot 4{,}25/8=1{,}859\cdot 10^{10}$~bytes.
  \item Bytes por paso en (i): $18{,}59+1{,}879=20{,}47$~GB; cota: $6{,}7\cdot 10^{12}/20{,}47\cdot 10^{9}=327{,}3$ tokens/s, $3{,}511$ veces más que en (i).
  \item Especulativa: $0{,}75^{6}=0{,}1780$; $\E[N_{\mathrm{r}}]=(1-0{,}1780)/0{,}25=3{,}288$ tokens por ronda; coste $5\cdot 0{,}04+1=1{,}2$ pasadas; aceleración $3{,}288/1{,}2=2{,}740$.
\end{enumerate}
\resultado (a) 70~GB de pesos y 80~GB libres; 224~KiB por token y $\approx 1{,}88$~GB ($1{,}75$~GiB) por secuencia. (b) 42 secuencias con GQA y solo 6 con MHA. (c) (i) $10{,}7$~ms por paso y $93{,}2$ tokens/s; (ii) $14{,}9$~ms por paso, $\approx 2140$ tokens/s totales y $\approx 67{,}0$ por usuario; el cómputo ($1{,}13$~ms) no limita. El lote crítico es $B\approx 295$, que no cabe: su caché ocuparía unos 277~GB y solo caben 85 secuencias. La decodificación de este servidor estará siempre limitada por la memoria. (d) $\TTFT\approx 0{,}283$~s. (e) $\approx 327$ tokens/s con int4 ($3{,}51$ veces más); la decodificación especulativa añade un factor $\approx 2{,}74$.
\comprobacion (a) $229\,376\cdot 8192=1{,}75\cdot 2^{30}$ exactamente. (b) $42/6=7=56/8$. (c) El lote se multiplica por 32 y el tiempo por paso solo por $14{,}94/10{,}73=1{,}39$: los pesos se leen una vez para todo el lote. (d) La intensidad del \emph{prefill} es $2N\cdot 4000/(2N)=4000$~FLOP/byte, muy por encima de $\Icrit\approx 295$: limitado por el cómputo, como se supuso. (e) Sin caché, int4 daría $70/18{,}59=3{,}77$ veces; la caché, que sigue en bf16 y pesa el $1{,}879/20{,}47=9{,}2\,\%$ de los bytes, rebaja la mejora a $3{,}51$.

[Criterios: (a) 2~p. (b) 2~p. (c) 1~p por cada caso y 1~p por las comprobaciones de cómputo y del lote crítico. (d) 1~p. (e) 1~p por int4 y 1~p por la decodificación especulativa.]

Repasar: M19, §2.1, §2.2, §2.3, §2.4, §3.4 y §4.3; M17, §2.2; M18, §2.3.
\end{solucion}

% ---------------------------------------------------------------------
\begin{solucion}{eje:extrapolar}
\datos $L(10^{20})=3{,}2$; $L(10^{22})=2{,}6$; $L(C)=1{,}7+kC^{-b}$. LSH con $(r,b)=(8,16)$ y $(5,20)$. Respuestas de $k=20$ tokens con $p\in\{0{,}9;\,0{,}95;\,0{,}98;\,0{,}995\}$.
\planteamiento (a) El cociente de los excesos sobre $1{,}7$ elimina $k$ (M16, §3.2). (b) Curva en S $P(s)=1-(1-s^{r})^{b}$ (M16, §2.5); el umbral exacto resuelve $1-(1-s^{r})^{b}=1/2$, es decir, $s_{1/2}=(1-2^{-1/b})^{1/r}$. (c) $\PP(\text{respuesta correcta})=p^{k}=e^{k\ln p}$ (M16, §6.2).
\pasos

\noindent\textbf{(a)}
\begin{enumerate}[label=\arabic*.]
  \item Excesos: $3{,}2-1{,}7=1{,}5$ y $2{,}6-1{,}7=0{,}9$; $0{,}9/1{,}5=100^{-b}=0{,}6$.
  \item $b=\ln(1{,}5/0{,}9)/\ln 100=0{,}5108/4{,}605=0{,}1109$.
  \item $L(10^{24})=1{,}7+0{,}9\cdot 100^{-b}=1{,}7+0{,}9\cdot 0{,}6=2{,}24$.
  \item Para pasar de un exceso de $0{,}9$ a $0{,}45$ hace falta un factor $F$ con $F^{-b}=0{,}5$: $F=2^{1/b}=2^{9{,}015}=517{,}4$, luego $C=517{,}4\cdot 10^{22}=5{,}174\cdot 10^{24}$.
\end{enumerate}

\noindent\textbf{(b)}
\begin{enumerate}[label=\arabic*.]
  \item $s=0{,}5$: $s^{8}=0{,}003906$; $(1-0{,}003906)^{16}=0{,}9393$; $P=0{,}0607$.
  \item $s=0{,}7$: $s^{8}=0{,}05765$; $0{,}94235^{16}=0{,}3867$; $P=0{,}6133$. (Con bases tan próximas a~1 conviene no redondear la base.)
  \item $s=0{,}9$: $s^{8}=0{,}4305$; $0{,}5695^{16}=1{,}225\cdot 10^{-4}$; $P=0{,}9999$.
  \item Umbral aproximado: $(1/16)^{1/8}=e^{-2{,}773/8}=0{,}7071$. Umbral exacto: $2^{-1/16}=0{,}9576$, $1-0{,}9576=0{,}04240$, $s_{1/2}=0{,}04240^{1/8}=0{,}6736$.
  \item Configuración $(5,20)$: aproximado $(1/20)^{1/5}=0{,}5493$; exacto $(1-2^{-1/20})^{1/5}=0{,}5087$.
\end{enumerate}

\noindent\textbf{(c)}
\begin{enumerate}[label=\arabic*.]
  \item $p=0{,}9$: $20\ln p=-2{,}107$; $p^{20}=0{,}1216$.
  \item $p=0{,}95$: $20\ln p=-1{,}026$; $p^{20}=0{,}3585$.
  \item $p=0{,}98$: $20\ln p=-0{,}4041$; $p^{20}=0{,}6676$.
  \item $p=0{,}995$: $20\ln p=-0{,}1003$; $p^{20}=0{,}9046$.
\end{enumerate}
\resultado (a) $b\approx 0{,}111$; $L(10^{24})\approx 2{,}24$; hace falta $C\approx 5{,}17\cdot 10^{24}$~FLOP, unas 517 veces el cómputo de $10^{22}$. Extrapolar cuatro órdenes de magnitud con dos puntos es arriesgado: el término irreducible $1{,}7$ se ha supuesto, no ajustado. (b) $P=0{,}0607$; $0{,}613$ y $0{,}9999$; umbral aproximado $0{,}707$ y exacto $0{,}674$, frente a $0{,}549$ y $0{,}509$ con $(5,20)$: la configuración $(8,16)$ es más estricta. Aumentan los falsos negativos (casi duplicados con $s\approx 0{,}6$ escapan más a menudo) y disminuyen los falsos positivos que hay que verificar. (c) $p^{20}=0{,}122$; $0{,}358$; $0{,}668$; $0{,}905$, y $20\ln p=-2{,}11$; $-1{,}03$; $-0{,}404$; $-0{,}100$. La log-probabilidad mejora de forma suave (es lineal en $\ln p$), mientras que la exactitud exacta permanece baja y después se dispara. Schaeffer, Miranda y Koyejo sostienen que muchas capacidades «emergentes» son un efecto de métricas discontinuas o no lineales como esta: con métricas continuas, la curva vuelve a ser suave.
\comprobacion (a) $L(5{,}174\cdot 10^{24})=1{,}7+0{,}9\cdot 517{,}4^{-0{,}1109}=1{,}7+0{,}9\cdot 0{,}5=2{,}15$. (b) En $s=(1/b)^{1/r}$, $s^{r}=1/16$ y $P=1-(15/16)^{16}=0{,}6439$, cerca del punto de máxima pendiente; el umbral exacto queda por debajo del aproximado, como en el ejercicio~3.4.2 de~\cite{mmds}. (c) $p^{20}=e^{20\ln p}$: las dos columnas contienen la misma información en escalas distintas; el valor para $p=0{,}95$ coincide con la cuestión~18.

[Criterios: (a) 2~p por $b$, 1~p por la predicción y 1~p por el cómputo; mencionar el riesgo de extrapolar suma, no resta. (b) 2~p por los cálculos y 1~p por la discusión. (c) 2~p por las tablas y 1~p por la interpretación.]

Repasar: M16, §2.5, §3.1, §3.2, §6.1 y §6.2.
\end{solucion}

% ---------------------------------------------------------------------
\section{Proyecto: valores de referencia para la corrección}

$C=2048\cdot 60\cdot 86\,400\cdot 989\cdot 10^{12}\cdot 0{,}4\approx 4{,}20\cdot 10^{24}$~FLOP; $2{,}95\cdot 10^{6}$ GPU-horas; 2477~MWh; $7{,}37\cdot 10^{6}$~\$. Óptimo de la ley de Chinchilla: $N\approx 7{,}90\cdot 10^{10}$ y $D\approx 8{,}87\cdot 10^{12}$. Regla de 20: $N\approx 1{,}87\cdot 10^{11}$ y $D\approx 3{,}74\cdot 10^{12}$. Ambos repartos caben en los $1{,}2\cdot 10^{13}$ tokens únicos. La pérdida prevista apenas cambia entre $N=5\cdot 10^{10}$ ($1{,}869$) y $N=10^{11}$ ($1{,}868$), de modo que un $N$ menor que el óptimo es defendible por el coste de inferencia. Cualquier elección bien argumentada es válida.
